\subsection{Passage from the local to the global}

PROPOSITION 8. Let A be a ring, m a maximal ideal of A and M an A-module. If there exists an ideal a of A such that m is the only maximal ideal of A containing a and aM = 0, then the canonical homomorphism \(M \rightarrow M_{m}\) is bijective.

A/a is then a local ring with maximal ideal m/a; M can be considered as an (A/a)-module; for all \(s \in A - m\) the canonical image of s in A/a is invertible, hence the homothety \(x \mapsto sx\) of M is bijective from the definition of \(M_{m}\) as the solution of a universal problem (§ 2, no. 2); whence the proposition.

In particular, if there exists \(k \geqslant 0\) such that \(\mathfrak{m}^k\mathbf{M} = 0\) , the homomorphism \(\mathbf{M} \to \mathbf{M}_m\) is bijective (§ 1, no. 1, Corollary to Proposition 1).

PROPOSITION 9. Let A be a ring, m a maximal ideal of A, M an A-module and k an integer \(\geqslant0\) . The canonical homomorphism \(M\to M_{m}/m^{k}M_{m}\) is surjective, has kernel \(m^{k}M\) and defines an isomorphism of \(M/m^{k}M\) onto \(M_{m}/m^{k}M_{m}\) .

Since the case \(k=0\) is trivial, suppose that \(k\geqslant1\) . It follows from Proposition 8 that the canonical homomorphism \(\mathbf{M}/\mathfrak{m}^{k}\mathbf{M}\to(\mathbf{M}/\mathfrak{m}^{k}\mathbf{M})_{\mathfrak{m}}\) is bijective. On the other hand \((\mathbf{M}/\mathfrak{m}^{k}\mathbf{M})_{\mathfrak{m}}\) is canonically identified with \(\mathbf{M}_{\mathfrak{m}}/(\mathfrak{m}^{k}\mathbf{M})_{\mathfrak{m}}\) (§2, no. 4, Theorem 1) and hence \((\mathfrak{m}^{k}\mathbf{M})_{\mathfrak{m}}=\mathfrak{m}^{k}\mathbf{M}_{\mathfrak{m}}\) (§2, no. 7, Corollary to Proposition 18), whence there is an isomorphism of \(\mathbf{M}/\mathfrak{m}^{k}\mathbf{M}\) onto \(\mathbf{M}_{\mathfrak{m}}/\mathfrak{m}^{k}\mathbf{M}_{\mathfrak{m}}\) which maps the class of an element \(x\in\mathbf{M}\) to the class of \(x/1\) .

COROLLARY. Let A be a ring, \(m_{1}, m_{2}, \ldots, m_{n}\) distinct maximal ideals of A, M an A-module and \(k_{1}, k_{2}, \ldots, k_{n}\) integers \(\geqslant 0\) . The canonical homomorphism from M to \(\prod_{i=1}^{n} \mathbf{M}_{\mathfrak{m}_{i}}/\mathfrak{m}_{i}^{k_{i}} \mathbf{M}_{\mathfrak{m}_{i}}\) is surjective and its kernel is \(\left(\bigcap_{i=1}^{n} \mathfrak{m}_{i}^{k_{i}}\right) \mathbf{M}\) .

This follows easily from Proposition 9 and §1, no. 2, Proposition 6, the \(\mathfrak{m}_{i}^{k_{i}}\) being relatively prime in pairs (§1, no. 2, Proposition 3).

In the rest of this no. A will denote a ring and \(\Omega(A)\) (or \(\Omega\) ) the set of maximal ideals of \(A\) .

PROPOSITION 10. The A-module \(\bigoplus_{m\in\Omega}A_m\) , the direct sum of the \(A_m\) , where \(m\in\Omega\) , is faithfully flat.

Each of the \(A_{m}\) is a flat A-module ( \(\S\) 2, no. 4, Theorem 1), hence \(E = \bigoplus_{m \in \Omega} A_{m}\) is flat (Chapter I, \(\S\) 2, no. 3, Proposition 2). Moreover, for every maximal ideal m of A, \(mA_{m}\) is the unique maximal ideal of \(A_{m}\) , hence \(mA_{m} \neq A_{m}\) , whence it follows that \(mE \neq E\) and consequently E is faithfully flat (Chapter I, \(\S\) 3, no. 1, Proposition 1 (d)).

THEOREM 1. Let M, N be two A-modules, \(u: \mathbf{M} \to \mathbf{N}\) an A-homomorphism and, for all \(m \in \Omega\) , let \(u_m: \mathbf{M}_m \to \mathbf{N}_m\) be the corresponding \(\mathbf{A}_m\) -homomorphism (§ 2, no. 2, Remark 5). For \(u\) to be injective (resp. surjective, bijective, zero), it is necessary and sufficient that, for all \(m \in \Omega\) , \(u_m\) be injective (resp. surjective, bijective, zero).

To say that, for all \(m \in \Omega\) , \(u_m\) is injective (resp. surjective, bijective, zero) is equivalent to saying that the homomorphism \(\bigoplus_{m} u_m: \bigoplus_{m} M_m \to \bigoplus_{m} N_m\) has the same property. But \(\bigoplus_{m} M_m = M \otimes_A E\) , \(\bigoplus_{m} N_m = N \otimes_A E\) and \(\bigoplus_{m} u_m = u \otimes 1\) , where \(E = \bigoplus_{m} A_m\) ; as \(E\) is faithfully flat (Proposition 10), the theorem follows from Chapter I, § 3, no. 1, Proposition 1 (c) and Proposition 2.

COROLLARY 1. Let M be an A-module, N a submodule of M and x an element of M. For \(x \in N\) , it is necessary and sufficient that, for all \(m \in \Omega\) , the canonical image of x in \(M_{m}\) belong to \(N_{m}\) .

Let \(\bar{x}\) be the class of \(x\) in M/N; to say that \(x \in N\) means that the A-linear mapping \(u: \alpha \mapsto \alpha \bar{x}\) from A to M/N is zero. Now, \((\mathrm{M} / \mathrm{N})_{\mathfrak{m}}\) is identified with \(\mathrm{M}_{\mathfrak{m}} / \mathrm{N}_{\mathfrak{m}}\) (§ 2, no. 4, Theorem 1) and \(u_{\mathfrak{m}}: \mathrm{A}_{\mathfrak{m}} \to \mathrm{M}_{\mathfrak{m}} / \mathrm{N}_{\mathfrak{m}}\) with the mapping \(\lambda \mapsto \lambda \bar{x}_{\mathfrak{m}}\) , where \(\bar{x}_{\mathfrak{m}}\) is the class mod. \(\mathrm{N}_{\mathfrak{m}}\) of the canonical image of \(x\) in \(\mathrm{M}_{\mathfrak{m}}\) . As the relation \(u = 0\) is equivalent to \(u_{\mathfrak{m}} = 0\) for all m by Theorem 1, this proves the corollary.

COROLLARY 2. Let \(\mathbf{M}\) be an A-module and, for all \(\mathfrak{m} \in \Omega\) , let \(f_{\mathfrak{m}}\) be the canonical mapping \(\mathbf{M} \to \mathbf{M}_{\mathfrak{m}}\) . The homomorphism \(x \mapsto (f_{\mathfrak{m}}(x))\) of \(\mathbf{M}\) to \(\prod_{\mathfrak{m} \in \Omega} \mathbf{M}_{\mathfrak{m}}\) is injective.

Applying Corollary 1 in the case N = 0, it is seen that the relation x = 0 is equivalent to \(f_{m}(x) = 0\) for all \(m \in \Omega\) .

COROLLARY 3. (i) Let \(\mathfrak{b}\) be an ideal of \(\mathbf{A}\) and a an element of \(\mathbf{A}\) . For \(a \in \mathfrak{b}\) , it is necessary and sufficient that, for all \(\mathfrak{m} \in \Omega\) , the canonical image of \(a\) in \(\mathbf{A}_{\mathfrak{m}}\) belong to \(\mathfrak{b}\mathbf{A}_{\mathfrak{m}}\) .

\begin{enumerate}
\def\labelenumi{(\roman{enumi})}
\setcounter{enumi}{1}
\tightlist
\item
  In particular, let \(b\) and \(c\) be two elements of \(A\) . For \(c\) to be a multiple of \(b\) , it is necessary and sufficient that, for all \(m \in \Omega\) , the canonical image of \(c\) in \(A_m\) be a multiple of that of \(b\) .
\end{enumerate}

As \(\mathfrak{bA}_{m} = \mathfrak{b}_{m}\) (§ 2, no. 7, Corollary to Proposition 18), (i) is a special case of Corollary 1; (ii) follows from (i) applied to the ideal Ab.

COROLLARY 4. Let A be an integral domain, K its field of fractions and M a torsion-free A-module such that M is canonically identified with a sub-A-module of \(\mathbf{K} \otimes_{\mathbf{A}} \mathbf{M}\) . Then, for all \(\mathfrak{m} \in \Omega\) , \(\mathbf{M}_{\mathfrak{m}}\) is canonically identified with a sub-A-module of \(\mathbf{K} \otimes_{\mathbf{A}} \mathbf{M}\) and \(\mathbf{M} = \bigcap_{\mathfrak{m} \in \Omega} \mathbf{M}_{\mathfrak{m}}\) .

As M is identified with a submodule of \(K \otimes_{A} M\) , \(M_{m}\) is identified with a sub- \(A_{m}\) -module of \((\mathbf{K} \otimes_{\mathbf{A}} \mathbf{M})_{\mathfrak{m}} = \mathbf{K}_{\mathfrak{m}} \otimes_{\mathbf{A}} \mathbf{M}\) (§2, no. 4, Theorem 1); as \(K_{m} = K\) , we see straightway that \(M_{m}\) is torsion-free; moreover, the commutativity of the diagram

\[
\begin{array}{c} \mathbf {M} \longrightarrow \mathbf {K} \otimes_ {\mathbf {A}} \mathbf {M} \\ \Big \downarrow \qquad \qquad \qquad \qquad \qquad \qquad \qquad \qquad \Big \downarrow \\ \mathbf {M} _ {\mathfrak {m}} \longrightarrow (\mathbf {K} \otimes_ {\mathbf {A}} \mathbf {M}) _ {\mathfrak {m}} \end{array}
\]

proves that the canonical mapping \(M \rightarrow M_{m}\) is injective. The corollary then follows from Corollary 1 applied to the A-module \(K \otimes_{A} M\) and its submodule M.

In particular, for every integral domain A,

\[
\mathbf {A} = \bigcap_ {\mathfrak {m} \in \Omega} \mathbf {A} _ {\mathfrak {m}}.\tag{2}
\]

COROLLARY 5. Let A be a ring. Every system of generators of the A-module \(\mathbf{A}^n\) with \(n\) elements is a basis of \(\mathbf{A}^n\) .

Let \((e_i)_{1\leqslant i\leqslant n}\) be the canonical basis of \(\mathbf{A}^n\) , \((x_i)_{1\leqslant i\leqslant n}\) a system of generators of \(\mathbf{A}^n\) with \(n\) elements and \(u\colon \mathbf{A}^n\to \mathbf{A}^n\) the A-linear mapping such that \(u(e_i) = x_i\) for \(1\leqslant i\leqslant n\) . By hypothesis \(u\) is surjective and it is necessary to show that \(u\) is injective. By virtue of Theorem 1 this can immediately be reduced to the case where \(\mathbf{A}\) is a local ring; if \(\mathfrak{m}\) is the maximal ideal of \(\mathbf{A}\) , the elements \(1\otimes x_i\) ( \(1\leqslant i\leqslant n\) ) in \((\mathrm{A / m})^n\) then form a system of generators of the free \((\mathrm{A / m})\) -module \((\mathrm{A / m})^n\) ; as \(\mathrm{A / m}\) is a field, this system is a basis of \((\mathrm{A / m})^n\) ; as \(\mathbf{A}^n\) is a free A-module, we deduce from Proposition 5 that \((x_i)\) is a basis of \(\mathbf{A}^n\) .

PROPOSITION 11. Let M be an A-module, N a finitely generated A-module and \(u: \mathbf{M} \to \mathbf{N}\) a homomorphism. For \(u\) to be surjective, it is necessary and sufficient that, for all \(\mathfrak{m} \in \Omega\) , the homomorphism \(\mathbf{M} / \mathfrak{m}\mathbf{M} \to \mathbf{N} / \mathfrak{m}\mathbf{N}\) derived from \(u\) by taking quotients be surjective.

It follows from Theorem 1 that, for \(u\) to be surjective, it is necessary and sufficient that \(u_{\mathfrak{m}}: \mathrm{M}_{\mathfrak{m}} \to \mathrm{N}_{\mathfrak{m}}\) be surjective for all \(\mathfrak{m} \in \Omega\) . As \(\mathrm{A}_{\mathfrak{m}}\) is a local ring and \(\mathrm{N}_{\mathfrak{m}}\) is a finitely generated \(\mathrm{A}_{\mathfrak{m}}\) -module, this amounts to saying that the homomorphism \(u_{\mathfrak{m}}': \mathrm{M}_{\mathfrak{m}} / \mathrm{mM}_{\mathfrak{m}} \to \mathrm{N}_{\mathfrak{m}} / \mathrm{mN}_{\mathfrak{m}}\) , obtained by taking quotients, is surjective (no. 2, Corollary 1 to Proposition 4); but \(\mathrm{M}_{\mathfrak{m}} / \mathrm{mM}_{\mathfrak{m}}\) (resp. \(\mathrm{N}_{\mathfrak{m}} / \mathrm{mN}_{\mathfrak{m}}\) ) is identified with \(\mathrm{M} / \mathrm{mM}\) (resp. \(\mathrm{N} / \mathrm{mN}\) ) (Proposition 9), whence the proposition.

PROPOSITION 12. Let E, F, G be three A-modules and \(v: \mathbf{G} \to \mathbf{F}\) , \(u: \mathbf{E} \to \mathbf{F}\) homomorphisms. Suppose that E is finitely presented. For there to exist a homomorphism \(w: \mathbf{E} \to \mathbf{G}\) such that \(u\) factors into \(\mathbf{E} \xrightarrow{w} \mathbf{G} \xrightarrow{v} \mathbf{F}\) , it is necessary and sufficient that, for all \(m \in \Omega\) , there exist a homomorphism \(w^m: \mathbf{E}_m \to \mathbf{G}_m\) such that \(u_m: \mathbf{E}_m \to \mathbf{F}_m\) factors into \(\mathbf{E}_m \xrightarrow{w^m} \mathbf{G}_m \xrightarrow{v_m} \mathbf{F}_m\) .

The existence of w satisfying the above statement is equivalent to the following property: u belongs to the image P of the mapping

\[
r = \operatorname{Hom} (1 _ {\mathbf {E}}, v): \operatorname{Hom} _ {\mathbf {A}} (\mathbf {E}, \mathbf {G}) \rightarrow \operatorname{Hom} _ {\mathbf {A}} (\mathbf {E}, \mathbf {F}).
\]

Now, \((\mathrm{Hom}_{\mathbf{A}}(\mathrm{E},\mathbf{F}))_{\mathfrak{m}}\) (resp. \((\mathrm{Hom}_{\mathbf{A}}(\mathrm{E},\mathbf{G}))_{\mathfrak{m}}\) ) is canonically identified with \(\mathrm{Hom}_{\mathbf{A}_{\mathfrak{m}}}(\mathrm{E}_{\mathfrak{m}},\mathrm{F}_{\mathfrak{m}})\) (resp. \(\mathrm{Hom}_{\mathbf{A}_{\mathfrak{m}}}(\mathrm{E}_{\mathfrak{m}},\mathrm{G}_{\mathfrak{m}})\) ) (§ 2, no. 7, Proposition 19 (i)), the canonical image of u in \((\mathrm{Hom}_{\mathbf{A}}(\mathrm{E},\mathbf{F}))_{\mathfrak{m}}\) is identified with \(u_{m}\) , \(r_{m}\) is identified with \(\mathrm{Hom}_{\mathbf{A}_{\mathfrak{m}}}(1_{\mathbf{E}_{\mathfrak{m}}},v_{\mathfrak{m}})\) and \(P_{m}\) with the image of \(r_{m}\) . The proposition then follows from Corollary 1 to Theorem 1 applied to \(\mathrm{Hom}_{\mathbf{A}}(\mathrm{E},\mathbf{F})\) and its submodule P.

COROLLARY 1. Let M be an A-module and N a submodule of M such that M/N is finitely presented. For N to be a direct factor of M, it is necessary and sufficient that, for all \(m \in \Omega\) , \(N_{m}\) be a direct factor of \(M_{m}\) .

To say that N is a direct factor of M means that the identity homomorphism on M/N factors into \(M/N \xrightarrow{w} M \xrightarrow{\phi} M/N\) where \(\phi\) is the canonical homomorphism and w a homomorphism (Algebra, Chapter II, § 1, no. 9, Proposition 14); as \((\mathrm{M}/\mathrm{N})_{\mathrm{m}} = \mathrm{M}_{\mathrm{m}}/\mathrm{N}_{\mathrm{m}}\) and \(\phi_{m}\) is the canonical homomorphism \(M_{m} \to M_{m}/N_{m}\) , the corollary follows easily from Proposition 12.

COROLLARY 2. Let M be a finitely generated free A-module and N a submodule of M which is a finitely generated free A-module. For N to be a direct factor of M, it is necessary and sufficient that, for all m ∈ Ω, mN = N ∩ (mM).

By definition M/N is finitely presented; on the other hand, \(N_{m}\) and \(M_{m}\) are finitely generated free \(A_{m}\) -modules. For \(N_{m}\) to be a direct factor of \(M_{m}\) , it is necessary and sufficient that the canonical mapping \(N_{m}/mN_{m} \rightarrow M_{m}/mM_{m}\) be injective (no. 2, Proposition 6); this is the same as saying that the canonical mapping \(N/mN \rightarrow M/mM\) must be injective (Proposition 9), and as its kernel is \((N \cap mM)/mN\) , this proves the corollary.

Proposition 12 (resp. its Corollary 1) will be applied in particular when A is Noetherian and E (resp. M/N) a finitely generated A-module (Chapter I, § 2, no. 8, Lemma 8).

